% grokking-transformer -- paper-style write-up.
%
% Build:  ./paper/build.sh        (regenerates figures, then runs tectonic)
% CI builds this on every push; see .github/workflows/ci.yml, job "paper".
%
% STATUS: outline. Every \outline{...} block below is a plan for prose that is
% not written yet, and the \outlinetrue switch on the next line makes them
% visible in the PDF (with a DRAFT banner) so an unfinished draft cannot be
% mistaken for a finished one. Setting \outlinefalse hides them.
%
% Figures come from paper/figures/*.pdf, which is NOT committed: it is
% regenerated from the committed logs by experiments/paper_figures.py before
% every build, so a figure here can never be staler than the data it plots.

\documentclass[11pt]{article}

\usepackage[letterpaper,margin=1in]{geometry}
\usepackage{amsmath}
\usepackage{amssymb}
\usepackage{graphicx}
\usepackage{booktabs}
\usepackage{xcolor}
\usepackage[numbers]{natbib}
\usepackage[colorlinks=true,allcolors=blue!55!black]{hyperref}

\graphicspath{{figures/}}

% ---- outline scaffolding (removed when the prose lands) --------------------
\newif\ifoutline
\outlinetrue

\newcommand{\outline}[1]{%
  \ifoutline
    \par\medskip\noindent
    \begingroup\small\color{blue!45!black}%
    \textsf{\textbf{[outline]}}~#1\par
    \endgroup\medskip
  \fi}

\newcommand{\draftbanner}{%
  \ifoutline
    \begin{center}\small\color{red!60!black}
    \textsf{DRAFT --- section skeleton and figures only; the prose marked
    \textbf{[outline]} is not written yet.}
    \end{center}
  \fi}
% ---------------------------------------------------------------------------

\title{Grokking in a one-layer transformer:\\
       what controls the delay, and what forms during it}
\author{porth-bot\thanks{Source, logs and figures:
        \url{https://github.com/porth-bot/grokking-transformer}}}
\date{\today}

\begin{document}
\maketitle
\draftbanner

\begin{abstract}
\outline{Written on Day 20 with the framing sections, once the results
sections have fixed which numbers are the headline ones. It has to say: a
one-layer decoder-only transformer trained full-batch on 30\% of
$(a+b) \bmod 97$ reaches 100\% train accuracy at step ${\sim}100$ and stays
near chance on test for ${\sim}1{,}300$ steps (median over five seeds) before
jumping to 100\%; what moves that delay (weight decay, data fraction, and
whether the two trade off); where in the network the norm pressure has to act;
and what the circuit looks like on either side of the jump, with the
mechanistic read-outs measured over five seeds rather than shown on one. It
must also say plainly that the phenomenon is a reproduction of
\citet{power2022grokking} and the Fourier circuit a reproduction of
\citet{nanda2023progress} --- see the note in Section~\ref{sec:related}.}
\end{abstract}

% ===========================================================================
\section{Introduction}
\label{sec:intro}

\outline{Day 20. The setup in one paragraph, the phenomenon in one figure
(Figure~\ref{fig:main}), and an explicit statement of what this write-up adds
over the README: nothing about the phenomenon itself, which is known, and
everything about (a) the controls --- a dropout arm, a weight-decay-scope arm,
a learning-rate arm, a modulus arm --- and (b) error bars on mechanistic
read-outs that are usually reported from a single run. Close with the
contributions list, phrased so that no item claims discovery.}

\begin{figure}[t]
  \centering
  \includegraphics[width=\linewidth]{grokking_main.pdf}
  \caption{The phenomenon, on the single run whose checkpoints the repo
  commits (seed~0, $(a+b) \bmod 97$, 30\% training data, weight decay $1.0$).
  \emph{Left:} train and test accuracy against step (log $x$), with the
  memorization and generalization steps marked. Training accuracy saturates
  early; test accuracy stays near chance for roughly an order of magnitude in
  step count and then jumps. \emph{Right:} the $\ell_2$ norm of all
  parameters over the same trajectory, rising to the transition and falling
  afterwards. Seed spread is not shown here --- it is in Figures~\ref{fig:wd}
  and~\ref{fig:frac}.}
  \label{fig:main}
\end{figure}

% ===========================================================================
\section{Setup}
\label{sec:setup}

\outline{Day 20. Model: a one-layer decoder-only transformer
\citep{vaswani2017attention}, $d_{\text{model}} = 128$, 4 heads, 223k
parameters, attention and LayerNorm written out by hand and checked against the
fused PyTorch reference. Data: all
$p^2$ pairs for $p = 97$, sequence $[a, b, =]$, supervised at the ``='' position,
30\% train split, noiseless and exhaustive so ``test accuracy'' is unambiguous.
Optimization: full-batch AdamW \citep{loshchilov2019adamw} at
lr $10^{-3}$ --- full batch specifically so minibatch noise is not a confound,
and decoupled decay specifically because L2-through-Adam is a different
regularizer. Define the grok step (first eval step at which test accuracy
crosses the threshold) and say what is done with runs that never cross it
(right-censored, not dropped and not imputed --- Section~\ref{sec:surface}).
State the seed protocol: five seeds per cell everywhere except the
Section~\ref{sec:surface} grid, which is three.}

\begin{figure}[t]
  \centering
  \includegraphics[width=\linewidth]{grokking_loss.pdf}
  \caption{The same seed-0 run in cross-entropy rather than accuracy (log
  both axes). Train loss falls steadily while test loss stays flat and then
  collapses at the transition, which is why the delay is legible as overfitting
  before it is legible as generalization.}
  \label{fig:loss}
\end{figure}

% ===========================================================================
\section{What controls the delay}
\label{sec:controls}

\subsection{Weight decay}
\label{sec:wd}

\outline{Day 19. README \S1. The wd sweep and what it establishes: whether
grokking happens at all, and when. Report medians with min--max over five
seeds. Carry over the Day-17 correction explicitly --- the claim ``$wd = 0$
never groks'' is a statement about the 30\% row, not about the method, and
Section~\ref{sec:surface} is where that is shown.}

\begin{figure}[t]
  \centering
  \includegraphics[width=\linewidth]{wd_sweep.pdf}
  \caption{Test accuracy against step at three weight decays, 30\% training
  data. Line is the median over five seeds, band the interquartile range; the
  legend gives each arm's grok step as median [min--max] over those seeds. The
  $wd = 0$ arm never reaches the threshold inside the budget and is reported as
  such rather than extrapolated.}
  \label{fig:wd}
\end{figure}

\subsection{Data fraction}
\label{sec:frac}

\outline{Day 19. README \S2: less data, longer delay, five seeds per cell.}

\begin{figure}[t]
  \centering
  \includegraphics[width=\linewidth]{frac_sweep.pdf}
  \caption{Data fraction at weight decay $1.0$. \emph{Left:} test-accuracy
  trajectories at 25/30/40\% training data, median over five seeds with an
  interquartile band. \emph{Right:} steps to 99\% test accuracy against
  fraction, median with min--max whiskers. The 60\% point is a single-seed
  context point at the end where the delay has essentially vanished, and
  carries no error bar for that reason.}
  \label{fig:frac}
\end{figure}

\subsection{Do they trade off? The $wd \times \mathrm{frac}$ surface}
\label{sec:surface}

\outline{Day 19, and this is the subsection that needs the most care because
its interesting half is a \emph{negative}. Two results, and they must not be
run together. (i) Data substitutes for weight decay: at $wd = 0$, fraction
0.25 and 0.30 never grok inside the 25k-step budget (3/3 seeds each) while
0.40 groks in all three. (ii) The apparent \emph{interaction} does not
survive its own error bars: a log-additive null fits the nine uncensored cells
to 1.24$\times$ rms against seed spreads of 1.15--16.5$\times$, the three
censored cells look like a refutation, and refitting at the corners of the
open cells' observed seed ranges gives envelopes that all straddle 1. Say so.
Also explain the censoring convention here rather than in a footnote:
non-grokking cells are right-censored at the budget, sorted to $+\infty$ for
the median, and hatched rather than coloured in the heatmap so no gradient can
be read across them.}

\begin{figure}[t]
  \centering
  \includegraphics[width=\linewidth]{wd_frac_surface.pdf}
  \caption{The $(wd, \text{fraction})$ grid, three seeds per cell.
  \emph{Left:} median grok step per cell with the min--max range, hatched
  where no seed generalized inside the 25k-step budget --- those cells are
  right-censored at the budget, not interpolated over and not imputed to it.
  \emph{Right:} the same medians as curves per fraction, where the
  log-additive null of Section~\ref{sec:surface} predicts parallel lines.
  Arrows mark censored cells as lower bounds.}
  \label{fig:surface}
\end{figure}

\subsection{Controls: learning rate, regularizer, modulus}
\label{sec:controls-arms}

\outline{Day 19. Three arms whose job is to rule things out rather than to
show something. README \S3 (a $10\times$ learning-rate change leaves the
phenomenon intact), \S6 (dropout as an alternative regularizer --- is it the
norm specifically or any regularizer?), and \S4 (the delay does not grow with
the modulus; data size is the variable that matters, and this is a table, not
a figure).}

\begin{figure}[t]
  \centering
  \includegraphics[width=\linewidth]{lr_sweep.pdf}
  \caption{Learning-rate sensitivity: test accuracy at lr
  $3\times10^{-4}$, $10^{-3}$ and $3\times10^{-3}$, a $10\times$ range. The
  phenomenon survives throughout and the grok step tracks the learning rate
  (5{,}500 / 1{,}700 / 800). One seed per arm.}
  \label{fig:lr}
\end{figure}

\begin{figure}[t]
  \centering
  \includegraphics[width=\linewidth]{dropout_control.pdf}
  \caption{The regularizer control. Dropout $0.1$ at $wd = 0$, against the
  main $wd = 1$ run and against the unregularized run. \emph{Left:} test
  accuracy --- dropout generalizes too. \emph{Right:} parameter norm (log),
  which is where the two regularizers part company: under dropout the norm
  grows rather than shrinks, so generalization here does not require the norm
  to come down.}
  \label{fig:dropout}
\end{figure}

\subsection{Where the norm pressure has to act}
\label{sec:scope}

\outline{Day 19. README \S7, and one of the results least covered elsewhere:
applying the decay to the embeddings only, versus everything but the
embeddings, versus all parameters. The readout circuit is where the pressure
has to be. Related to \citet{liu2023omnigrok}'s norm account and worth citing
against it.}

\begin{figure}[t]
  \centering
  \includegraphics[width=\linewidth]{wd_scope.pdf}
  \caption{Scope of the decay, at 30\% data and $wd = 1.0$, seed~0: applied
  to all parameters, to the embeddings only, and to everything except the
  embeddings. \emph{Left:} test accuracy. \emph{Right:} total parameter
  norm.}
  \label{fig:scope}
\end{figure}

% ===========================================================================
\section{What forms during the delay}
\label{sec:mechanism}

\subsection{Norm and Fourier structure}
\label{sec:fourier}

\outline{Day 19. README \S5. The weight-norm trajectory and the emergence of a
sparse set of frequencies in the embeddings, memorization checkpoint versus
final checkpoint. Carry the Day-17 correction: the rise-then-fall in weight
norm is a \emph{weight-decay} signature, not a grokking one --- the $wd = 0$
runs that do generalize show 0.00\% maximum drawdown.}

\begin{figure}[t]
  \centering
  \includegraphics[width=\linewidth]{fourier_spectrum.pdf}
  \caption{Discrete Fourier spectrum of the token embeddings at the
  memorization checkpoint (left) and at the end of training (right), seed~0.
  The spectrum goes from near-flat to concentrated on a handful of
  frequencies; the panels are annotated with the share of energy carried by
  the top five.}
  \label{fig:spectrum}
\end{figure}

\begin{figure}[t]
  \centering
  \includegraphics[width=\linewidth]{embedding_circle.pdf}
  \caption{Token embeddings projected onto the $\cos$/$\sin$ plane of the
  dominant frequency ($k = 36$), coloured by the residue $kn \bmod p$, before
  and after generalization (seed~0). Panels are annotated with the in-plane
  variance share and the radial coefficient of variation, which are the two
  numbers that make ``ring'' a measurement rather than a description.}
  \label{fig:circle}
\end{figure}

\subsection{The generalizing solution is a sparse sum over frequencies}
\label{sec:attribution}

\outline{Day 19. README \S8, the logit-attribution result and the one number
from it that matters most for the paper's argument: restricting the
\emph{memorizing} model's logits to the $a+b$ subspace already recovers
accuracy well above what the raw model scores, so the circuit is forming
underneath the memorization before the test accuracy moves. Seed-0 values are
in the figure; the five-seed values are in Section~\ref{sec:seeds} and are the
ones to quote.}

\begin{figure}[t]
  \centering
  \includegraphics[width=\linewidth]{logit_attribution.pdf}
  \caption{Logit attribution at the two checkpoints of seed~0.
  \emph{Left:} share of the $a+b$ diagonal's logit energy per frequency.
  \emph{Right:} test accuracy of the model with its logits restricted to the
  top-$k$ $a+b$ frequencies, against $k$. The memorizing model is already well
  above its own unrestricted accuracy under that restriction, which is the
  form the ``circuit forms underneath the memorization'' claim actually
  takes.}
  \label{fig:attribution}
\end{figure}

\subsection{Progress measures}
\label{sec:progress}

\outline{Day 19. README \S10: the circuit forms gradually and the read-outs
move before the accuracy does. Relate to \citet{nanda2023progress} explicitly,
including where this reproduction agrees with theirs and where the setup
differs.}

\begin{figure}[t]
  \centering
  \includegraphics[width=\linewidth]{progress_measures.pdf}
  \caption{Progress measures against step, seed~0, with the grok step
  marked. \emph{Left:} full-model test loss against the loss of the model
  restricted to five $a+b$ frequencies and the loss with those frequencies
  removed. \emph{Right:} top-5 embedding energy fraction against test
  accuracy on a shared $x$ axis. Both read-outs move well before the accuracy
  does.}
  \label{fig:progress}
\end{figure}

\subsection{Attention}
\label{sec:attention}

\outline{Day 19. README appendix. The attention pattern before and after, and
the entropy trajectory. Keep this short: Section~\ref{sec:seeds} finds the
attention read-outs converge to the same place in every arm, so their role in
the paper is as a null, not as a finding.}

\begin{figure}[t]
  \centering
  \includegraphics[width=\linewidth]{attention_pattern.pdf}
  \caption{Attention weight from the ``='' position to each of $a$, $b$ and
  ``='', per head, at the memorization checkpoint and at the end of training
  (seed~0). Panels are annotated with the total operand weight and the $a$/$b$
  asymmetry.}
  \label{fig:attnpattern}
\end{figure}

\begin{figure}[t]
  \centering
  \includegraphics[width=\linewidth]{attention_entropy.pdf}
  \caption{The attention read-out along the trajectory (seed~0; $x$ is
  linear below step 100 and logarithmic above). \emph{Left:} entropy of the
  ``='' row, full and restricted to the operands, against the uniform and
  symmetric reference levels. \emph{Right:} operand weight and per-head
  $a$/$b$ asymmetry against test accuracy.}
  \label{fig:attnentropy}
\end{figure}

\subsection{Are the mechanistic numbers typical?}
\label{sec:seeds}

\outline{Day 19, and this is the section that earns the paper the right to
quote the rest. README \S12: every mechanistic read-out re-measured over five
seeds. Twelve of thirteen separate memorization from generalization completely
at $p = 0.008$, which is the \emph{floor} for five versus five --- so the
rank-sum helper reports $\min p = 1/\binom{n+m}{n}$ beside the $p$-value and
``not significant'' can be told from ``cannot be significant at this sample
size''. Then the honest half: on 8 of 13 read-outs seed~0 held the extreme in
the flattering direction, and the cause is a budget mismatch, not a lucky seed
--- seed~0's run was extended to 11,100 steps where the others early-stop at
1,800--2,100, and the seven read-outs where it is an outlier are the seven that
further consolidation would move.}

\begin{figure}[t]
  \centering
  \includegraphics[width=\linewidth]{mechanistic_seeds.pdf}
  \caption{Every mechanistic read-out at the memorization and final
  checkpoints over five seeds, with the exact rank-sum $p$-value per read-out
  ($0.008$ is the floor at five versus five). Bars are medians, dots are
  seeds, and the star marks seed~0 --- the published run, which was trained to
  11.1k steps where the other four early-stop near 2k, and which sits at the
  flattering extreme on most of these panels for that reason.}
  \label{fig:seeds}
\end{figure}

% ===========================================================================
\section{Scope of the result}
\label{sec:scope-of-result}

\subsection{Head count}
\label{sec:heads}

\outline{Day 19. README \S9: one wide head groks sooner than four narrow ones,
medians 300/700/1300 for 1/2/4 heads at five seeds per arm with complete
separation ($p = 0.008$, the floor). Two points to make and neither is
optional: the ordering separates because the within-arm spread
(1.3--1.6$\times$) is smaller than the between-arm step (2.3$\times$,
1.9$\times$), which is readable off the table without any test; and the
\emph{mechanistic} guess that went with this result did not survive --- all
four attention read-outs reach the same place in every arm and separate no
pair (smallest $p = 0.056$), so head count changes when the circuit forms, not
what it converges to.}

\begin{figure}[t]
  \centering
  \includegraphics[width=\linewidth]{head_count.pdf}
  \caption{Head count at fixed $d_{\text{model}} = 128$, so head width
  falls 128/64/32 as the count rises, five seeds per arm. \emph{Left:} median
  test accuracy with an interquartile band. \emph{Middle:} grok step per
  seed, bar at the median --- the arms separate completely. \emph{Right:}
  operand asymmetry of the grokked model, which does \emph{not} separate
  them.}
  \label{fig:heads}
\end{figure}

\subsection{Other operations}
\label{sec:ops}

\outline{Day 19. README \S11: subtraction and multiplication against addition,
and the swap-equivariance control that asks whether what is learned is the
group structure rather than commutativity. The subtraction arm is partly
censored and must be reported as such.}

\begin{figure}[t]
  \centering
  \includegraphics[width=\linewidth]{operations.pdf}
  \caption{Addition, subtraction and multiplication mod 97 at 30\% data and
  two weight decays, median of three seeds with a min--max band. Subtraction
  is the slow arm in both panels and is partly censored at $wd = 0.1$.}
  \label{fig:ops}
\end{figure}

\begin{figure}[t]
  \centering
  \includegraphics[width=\linewidth]{swap_equivariance.pdf}
  \caption{Operand-swap read-outs for addition and subtraction, at
  memorization and after grokking ($wd = 1$, three seeds, dots are seeds),
  normalized by the logit sd and against a shuffle baseline.
  \emph{Left:} invariance, $L(a,b)$ against $L(b,a)$ --- addition acquires it.
  \emph{Right:} anti-equivariance, $L(a,b)$ against $L(b,a)(-c)$ ---
  subtraction acquires that instead. So what is learned is the symmetry the
  operation has, not commutativity.}
  \label{fig:swap}
\end{figure}

% ===========================================================================
\section{Related work}
\label{sec:related}

\outline{Day 20, and the section with the highest risk of overclaiming, so it
gets written before the introduction rather than after. It must say without
hedging: the phenomenon is \citet{power2022grokking}'s, the Fourier circuit
and the progress-measure framing are \citet{nanda2023progress}'s, the norm
account is \citet{liu2023omnigrok}'s and the efficiency account is
\citet{varma2023efficiency}'s. This work reproduces them from a hand-written
implementation and adds controls and error bars; nothing here is a discovery
and the paper must not read as if it were.}

% ===========================================================================
\section{Limitations}
\label{sec:limitations}

\outline{Day 20. Named, not buried. One operation family on one modulus at one
model size; a 25k-step budget that right-censors part of the
Section~\ref{sec:surface} grid and is itself the reason the ``cliff'' there is
not a dynamical claim; three seeds rather than five on that grid; five seeds
being enough to reach $p = 0.008$ and not enough to reach anything smaller;
mechanistic read-outs taken at two checkpoints rather than continuously; and
the seed-0 budget mismatch of Section~\ref{sec:seeds}, which is fixed in the
reporting but not in the underlying runs.}

% ===========================================================================
\section{Conclusion}
\label{sec:conclusion}

\outline{Day 21, last, once everything above is settled and only then.}

% ===========================================================================
\bibliographystyle{plainnat}
\bibliography{refs}

\end{document}
